% ------------------------------------------------------------------
%  Preambule : packages, mise en page et commandes du rapport
% ------------------------------------------------------------------
\usepackage[T1]{fontenc}
\usepackage[utf8]{inputenc}
\usepackage{lmodern}
\usepackage[french]{babel}
\usepackage[a4paper,margin=2cm,headheight=14pt]{geometry}
\usepackage{microtype}
\usepackage{amsmath,amssymb}
\usepackage{graphicx}
\usepackage{booktabs,tabularx,array,multirow}
\usepackage[table]{xcolor}
\usepackage{float}
\usepackage{placeins}
\usepackage{caption}
\usepackage{subcaption}
\usepackage{enumitem}
\usepackage[version=4]{mhchem}
\usepackage{siunitx}
\usepackage{csquotes}
\usepackage{titlesec}
\usepackage{titletoc}
\usepackage{fancyhdr}
\usepackage{tikz}
\usepackage{pgfplots}
\usepackage{hyperref}

% --- Francais : "Tableau" et legendes sans petites capitales ----------
\frenchsetup{SmallCapsFigTabCaptions=false}
\addto\captionsfrench{\renewcommand{\tablename}{Tableau}}

% --- Unites et nombres (virgule decimale) ---------------------------
\sisetup{
  output-decimal-marker = {,},
  group-separator       = {\,},
  per-mode              = symbol,
  range-phrase          = {~à~},
  list-final-separator  = {~et~},
  list-pair-separator   = {~et~},
  range-units           = single,
}
\DeclareSIUnit{\tour}{tr}

% --- TikZ / pgfplots ------------------------------------------------
\usetikzlibrary{arrows.meta,calc,positioning,babel}
\pgfplotsset{compat=1.18}
\pgfkeys{/pgf/number format/.cd,use comma,1000 sep={}}

% Couleurs communes aux figures
\definecolor{c2000}{HTML}{1F5FA8}
\definecolor{c3000}{HTML}{D4762C}
\definecolor{c4000}{HTML}{2E8B57}
\definecolor{c5000}{HTML}{8A4FA8}
\definecolor{substrat}{HTML}{B9BEC6}
\definecolor{verre}{HTML}{CFE3EA}
\definecolor{nlof}{HTML}{E07A7A}
\definecolor{lor}{HTML}{F2CF6B}
\definecolor{resinepos}{HTML}{D9534F}
\definecolor{resineexpo}{HTML}{F6B3B3}
\definecolor{metal}{HTML}{5A6B80}
\definecolor{alu}{HTML}{8C9BB0}

% --- Legendes, numerotation par sujet --------------------------------
\captionsetup{font=small,labelfont=bf,labelsep=endash,skip=4pt}
\captionsetup[table]{position=top}
\captionsetup[sub]{font=small,labelfont=bf}
\counterwithin{figure}{section}
\counterwithin{table}{section}
\setlength{\textfloatsep}{10pt plus 2pt minus 2pt}
\setlength{\floatsep}{8pt plus 2pt minus 2pt}
\setlength{\intextsep}{8pt plus 2pt minus 2pt}
\renewcommand{\topfraction}{0.9}
\renewcommand{\bottomfraction}{0.8}
\renewcommand{\textfraction}{0.07}
\renewcommand{\floatpagefraction}{0.8}

% --- Titres : chaque section est un sujet ("Sujet 1 : ...") ----------
\titleformat{\section}{\normalfont\Large\bfseries}{Sujet~\thesection~:}{0.4em}{}
\titlespacing*{\section}{0pt}{0pt}{10pt}
\titlespacing*{\subsection}{0pt}{10pt plus 2pt minus 2pt}{4pt}
\titlespacing*{\paragraph}{0pt}{6pt plus 2pt minus 1pt}{0.8em}
\titlecontents{section}[0em]{\addvspace{0.5em}\bfseries}
  {Sujet~\thecontentslabel~:~}{}{\hfill\contentspage}

% --- En-tete / pied de page -----------------------------------------
\pagestyle{fancy}
\fancyhf{}
\fancyhead[L]{\small Rapport TP Micro-332}
\fancyhead[R]{\small \nouppercase{\leftmark}}
\fancyfoot[C]{\small \thepage}
\renewcommand{\headrulewidth}{0.4pt}
\renewcommand{\sectionmark}[1]{\markboth{Sujet~\thesection~: #1}{}}

\hypersetup{
  colorlinks = true,
  linkcolor  = black,
  citecolor  = black,
  urlcolor   = blue!60!black,
  pdftitle   = {Rapport TP Micro-332},
  pdfauthor  = {Groupe 8},
}

\setlist{itemsep=1pt,topsep=3pt,parsep=0pt}
\setlength{\parskip}{3pt}

% --- Marqueurs de travail (rouges) -----------------------------------
\newif\ifversionfinale
\versionfinalefalse
% \acompleter{texte} -> [À compléter : texte] ; \acompleter{} -> [À compléter]
\newcommand{\acompleter}[1]{%
  \ifversionfinale\else{\color{red}\textbf{[À compléter%
    \if\relax\detokenize{#1}\relax\else{}~: #1\fi]}}\fi}
\newcommand{\averifier}[1]{%
  \ifversionfinale\else{\color{red}\textbf{[À vérifier%
    \if\relax\detokenize{#1}\relax\else{}~: #1\fi]}}\fi}

% --- Photos : s'affichent si le fichier existe, sinon un cadre rouge --
%  Deposer les images (jpg/png) dans photos/ avec le nom indique.
%  \photobox{chemin}                         : l'image seule
%  \photo{chemin}{legende}{label}            : une photo
%  \photopaire{ch1}{leg1}{lab1}{ch2}{leg2}{lab2} : deux photos cote a cote
\newcommand{\photobox}[1]{%
  \IfFileExists{#1}{%
    \includegraphics[width=\linewidth,height=5.5cm,keepaspectratio]{#1}%
  }{%
    \fbox{\parbox[c][3cm][c]{0.92\linewidth}{\centering\color{red}\small
      Photo à insérer :\\ \texttt{\detokenize{#1}}}}%
  }}
\newcommand{\photo}[3]{%
  \begin{figure}[htbp]
    \centering
    \begin{minipage}{0.6\linewidth}\centering\photobox{#1}\end{minipage}
    \caption{#2}\label{#3}
  \end{figure}}
\newcommand{\photopaire}[6]{%
  \begin{figure}[htbp]
    \centering
    \begin{minipage}[t]{0.48\linewidth}\centering\photobox{#1}
      \caption{#2}\label{#3}\end{minipage}\hfill
    \begin{minipage}[t]{0.48\linewidth}\centering\photobox{#4}
      \caption{#5}\label{#6}\end{minipage}
  \end{figure}}

% Raccourcis
\newcommand{\Dz}{\ensuremath{D_{0}}}
\newcommand{\Dc}{\ensuremath{D_{100}}}
\newcommand{\ez}{\ensuremath{e_{0}}}
